\subsection{Функциональное тестирование в контейнерах}

Перед нагрузочными экспериментами корректность модуля проверена локально в
контейнеризованном окружении. Окружение разворачивается через Docker Compose и
включает: контейнер с Apache Flink (менеджер заданий и менеджер задач), контейнер с
объектным хранилищем, совместимым с протоколом S3 (MinIO), и контейнер-клиент для
запуска заданий. Артефакт Paimon с модулем \texttt{paimon-vortex} собирается из
исходных текстов; нативная библиотека Vortex кросс-компилируется под целевую
платформу контейнера.

Проверялись следующие свойства.

\begin{enumerate}
    \item \textbf{Запись и чтение таблицы только для добавления.} Создание таблицы с
        форматом Vortex, загрузка данных, чтение, сверка числа строк и контрольных
        агрегатов с эталоном.
    \item \textbf{Таблица с первичным ключом и слияние версий.} Начальная загрузка,
        затем серия обновлений по случайным ключам с возрастающим полем
        последовательности; чтение должно вернуть для каждого ключа последнюю версию.
        Корректность сверялась с результатом той же последовательности операций на
        формате Parquet.
    \item \textbf{Проекция столбцов и проталкивание предикатов.} Запросы с выборкой
        части столбцов и с селективными предикатами разных типов (целочисленные
        диапазоны, равенство строк, проверка на null); сверка результатов с эталоном
        и проверка по журналам, что предикат действительно протолкнут.
    \item \textbf{Эволюция числа бакетов и слияние.} Изменение числа бакетов с
        перераспределением данных; принудительное слияние; чтение после слияния.
    \item \textbf{Граничные случаи кодеков.} Столбцы с высокой и низкой
        кардинальностью, с большой долей null, со строками-префиксами, с числами с
        плавающей точкой --- проверка, что кодеки Vortex (FastLanes, ALP, FSST)
        активируются и данные восстанавливаются без потерь.
\end{enumerate}

В ходе тестирования выявлен и устранён ряд проблем сборки и конфигурации (несовместимость
версий Java и Flink на клиентском процессе, идемпотентность повторных запусков,
требования к памяти вне кучи JVM для нативных буферов Vortex, особенности
кросс-компиляции нативной библиотеки под версию системной библиотеки C целевой
платформы); они подробно рассмотрены в подразделе~\ref{subsec:issues}. Также
установлено существенное ограничение: целевой файловой системой для нативной
библиотеки Vortex является объектное хранилище, совместимое с S3; распределённая
файловая система Hadoop (HDFS) и схема URI \texttt{hdfs://} не поддерживаются, ---
что определило выбор объектного хранилища в качестве бэкенда и в контейнерном
окружении, и на кластере. После устранения проблем все перечисленные проверки
пройдены: артефакт признан функционально корректным и пригодным для нагрузочного
тестирования.
